\documentclass[../cours_spe_trigo.tex]{subfiles}
\begin{document}
\setcounter{chapter}{4}
\chapter{Les valeurs remarquables : les reconstruire sans les apprendre aveuglément}
\section{Pourquoi $\cos(45^\circ)=\sin(45^\circ)=\sqrt2/2$ ?}
Construisons un triangle rectangle isocèle dont les deux côtés de l'angle droit valent $1$. Par Pythagore, l'hypoténuse vaut $\sqrt{1^2+1^2}=\sqrt2$. Ses angles aigus sont égaux et leur somme vaut $90^\circ$ ; chacun mesure $45^\circ=\pi/4$. Donc
\[
\boxed{\cos\frac\pi4=\sin\frac\pi4=\frac1{\sqrt2}=\frac{\sqrt2}{2}.}
\]
\textbf{Pourquoi rationaliser ?} Multiplier $1/\sqrt2$ par $\sqrt2/\sqrt2$ donne $\sqrt2/2$, écriture traditionnellement retenue.

\section{Pourquoi les nombres $1/2$ et $\sqrt3/2$ ?}
Dans un triangle équilatéral de côté $2$, la hauteur partage la base en deux segments de longueur $1$ et forme deux triangles rectangles. Par Pythagore, la hauteur vaut $\sqrt{2^2-1^2}=\sqrt3$. Les angles de chacun sont $30^\circ$ et $60^\circ$. Par conséquent :
\[
\boxed{\begin{aligned}
\cos\frac\pi6&=\frac{\sqrt3}{2},&\sin\frac\pi6&=\frac12,\\
\cos\frac\pi3&=\frac12,&\sin\frac\pi3&=\frac{\sqrt3}{2}.
\end{aligned}}
\]
\begin{center}
\begin{tikzpicture}[scale=1]
\coordinate (A) at (0,0);\coordinate (B) at (4,0);\coordinate (C) at (2,3.464);
\draw[thick,bleuNuit] (A)--(B)--(C)--cycle;
\draw[densely dashed,orange,thick] (C)--(2,0);
\draw (2,0)--(2,.27)--(2.27,.27)--(2.27,0);
\node[below] at (1,0) {$1$};\node[below] at (3,0) {$1$};
\node[left] at (.9,1.82) {$2$};\node[right] at (3.1,1.82) {$2$};
\node[right,orange] at (2,1.73) {$\sqrt3$};
\node[above] at (2,3.55) {$30^\circ\;|\;30^\circ$};
\node[above right] at (.55,.15) {$60^\circ$};
\end{tikzpicture}
\end{center}
\section{Tableau construit et signes dans les quatre quadrants}
\begin{center}
\renewcommand{\arraystretch}{1.45}
\begin{tabular}{c|ccccc}
\toprule
$\theta$ & $0$ & $\pi/6$ & $\pi/4$ & $\pi/3$ & $\pi/2$ \\
\midrule
$\cos\theta$ & $1$ & $\sqrt3/2$ & $\sqrt2/2$ & $1/2$ & $0$\\
$\sin\theta$ & $0$ & $1/2$ & $\sqrt2/2$ & $\sqrt3/2$ & $1$\\
\bottomrule
\end{tabular}
\end{center}
Toutes les autres valeurs remarquables s'en déduisent par les symétries du cercle, sans nouveau triangle.
\end{document}
